\documentclass[english,a4paper]{my_memo2}
\synctex=1
\usepackage{mathtools}
\usepackage{microtype}
\usepackage{booktabs,longtable,array}
\usepackage{enumitem}
\hypersetup{hidelinks}

\newcommand{\nelec}{n_{\rm e}}
\newcommand{\nH}{n_{\rm H}}
\newcommand{\tcool}{t_{\rm cool}}
\newcommand{\Da}{\mathrm{Da}}

\title{\textbf{Radiative Turbulent Mixing Layers in Astrophysics}\\
\large Literature Review and a Practical Framework for Computing Emission Spectra}
\author{Kwang-Il Seon}
\date{\docmoddate}

\begin{document}
\maketitle

\begin{abstract}
Radiative turbulent mixing layers (RTMLs) occur where shearing hot and cold
gas mix, pass through the efficient-cooling temperature range, and radiate the
enthalpy supplied by the hot phase.  This note separates two related but
distinct calculations: a bolometric cooling loss and an observable spectrum.
It reviews the physical controls on RTMLs and presents an implementable
post-processing workflow: obtain the time-dependent thermodynamic history,
solve non-equilibrium ionization (NEI) where required, evaluate line and
continuum emissivities from an atomic database, integrate them over volume or
ray, and finally apply transfer and instrumental response.  The central
practical result is that a cooling table alone cannot predict a spectrum: the
ion fractions and level populations must be specified.  In rapidly mixed or
rapidly cooled gas, the ionization/recombination time can exceed the cooling
or mixing time, making CIE spectra a useful baseline but not the default
prediction.
For a new calculation, the most direct publicly documented implementation is
the full-ionization-network extension of \textsc{Pluto}; \textsc{Flash} is a
robust alternative but requires a new cooling/emissivity coupling, while
\textsc{Maihem} is the most integrated option if its modified-\textsc{Flash}
source can be obtained.  \textsc{Torus} is better retained for transfer and
synthetic-observation post-processing than used as the primary RTML
microphysics solver.  The locally available \textsc{Mappings~V} 5.2.1 source
substantially improves the practical options: it is not a multidimensional
mixing solver, but it can supply internally consistent CIE/NEI cooling
histories, ion fractions, and line/continuum emissivities, as well as
one-zone and radiative-shock benchmarks.  The most realistic modular
architecture is therefore \textsc{Pluto} or \textsc{Flash} for the RTML
dynamics, \textsc{Mappings} for microphysics and spectral tables, and
\textsc{Torus} only when a subsequent transfer calculation is required.  For
galactic-wind applications, the Fielding--Bryan steady multiphase-wind model
provides a useful additional closure between local interface emission and the
radial population of clouds, allowing a local RTML surface flux to be
converted into a galaxy-scale luminosity distribution.
\end{abstract}

\tableofcontents

\section{Scope and physical picture}

Consider cold gas $(T_c\sim10^4\,\mathrm{K})$ moving relative to a hot
phase $(T_h\gtrsim10^6\,\mathrm{K})$.  Kelvin--Helmholtz-driven turbulence
mixes material across the interface.  At nearly common pressure the mixed gas
passes through roughly
\begin{equation}
 T_{\rm mix}\sim (T_cT_h)^{1/2},
\end{equation}
often near $10^5\,\mathrm{K}$, where metal-line cooling is efficient.  In a
steady radiative layer, cooling is powered principally by the enthalpy flux of
entrained hot gas, with turbulent dissipation potentially contributing.  The
global energy statement is more robust than any particular geometric model:
\begin{equation}
 L_{\rm rad}=\int_V \nelec\nH\Lambda(T,\boldsymbol{x})\,dV
 \simeq \dot M_{\rm in}\,h_h+\dot E_{\rm diss}-\dot E_{\rm stored}.
 \label{eq:energy}
\end{equation}
Here $\Lambda$ is the \emph{net} cooling coefficient only when the assumed
ionization state and radiation field match the gas.  Equation~(\ref{eq:energy})
checks the energy budget; it does not determine how $L_{\rm rad}$ is distributed
over wavelength.

The relevant competition is captured by
\begin{equation}
 \Da\equiv\frac{t_{\rm turb}}{\tcool},\qquad
 \tcool=\frac{e_{\rm th}}{\nelec\nH\Lambda_{\rm net}}.
 \label{eq:da}
\end{equation}
For $\Da\gtrsim1$, cooling is fast compared with outer-scale mixing and the
interface fragments into a strongly corrugated multiphase cooling front.  Tan,
Oh, and Gronke emphasize that the outer eddy turnover time then sets the
effective consumption/mixing rate, whereas Fielding et al. find a corrugated
cooling surface with an approximately $5/2$ fractal dimension.  These are
physical guides for convergence tests, not replacements for a spectrum
calculation \cite{Tan2021,Fielding2020}.

\section{Development of the existing research}

\subsection{Analytic foundations and the first spectral predictions}

Begelman and Fabian proposed that shear-driven turbulent transport mixes hot
and cold material into an intermediate phase with a characteristic temperature
near $(T_cT_h)^{1/2}$, and that the incoming hot-phase enthalpy supplies the
radiative loss \cite{Begelman1990}.  This model established the two ideas that
remain central today: the interface is an energy-conversion layer rather than
an isolated parcel, and its luminosity is tied to the rate at which hot gas is
entrained.  Its assumed layer thickness, often written
$h\sim v_{\rm turb}\tcool$, should now be regarded as an early closure rather
than a universal scaling.

Slavin, Shull, and Begelman converted this picture into emission and absorption
predictions for Galactic interfaces \cite{Slavin1993}.  Their calculation
included time-dependent ionization and predicted C\,IV, Si\,IV, N\,V, O\,VI,
UV semiforbidden lines, and H$\alpha$.  This work was conceptually important
because it treated ionization lag and spectra together.  It also anticipated a
persistent observational issue: one individual interface generally produces a
modest high-ion column, so a long sightline may need to intersect many
interfaces.

\subsection{The numerical transition: dimensionality, NEI, and magnetic fields}

Esquivel et al. performed three-dimensional MHD calculations with equilibrium
cooling and demonstrated that field orientation and finite evolution time
matter for the morphology and growth of a mixing layer \cite{Esquivel2006}.
Kwak and Shelton subsequently coupled hydrodynamics to NEI rate equations
\cite{Kwak2010}.  They found that mixing can transport ions into gas whose
temperature differs from the CIE formation temperature, while cooling can
outpace recombination.  Their simulated C\,IV/N\,V/O\,VI ratios were broadly
compatible with Galactic-halo diagnostics, but varied strongly in space and
time.  Thus a single static ``mixing temperature'' is insufficient for either
ion columns or line emission.

Ji, Oh, and Masterson carried out three-dimensional HD and MHD simulations
including NEI and photoionization \cite{Ji2019}.  Their results revised several
early analytic expectations:
\begin{itemize}[leftmargin=*]
\item characteristic turbulent and inflow velocities were substantially below
the imposed shear velocity and only weakly dependent on it;
\item the layer width followed approximately $h\propto\tcool^{1/2}$ in their
parameter regime, rather than $h\propto\tcool$;
\item ion columns retained density and metallicity dependence; and
\item magnetic amplification and tension could strongly suppress mixing and
reduce high-ion columns.
\end{itemize}
NEI increased the abundances and columns of high ions by factors of a few in
relevant cases.  Photoionization can also over-ionize the gas and modify its
cooling efficiency, although Ji et al. found this secondary for the CGM
conditions they explored.

\subsection{Fractal fronts, combustion theory, and two cooling regimes}

Fielding et al. showed that most cooling occurs in a thin, corrugated sheet
whose measured area is consistent with a fractal dimension $D\simeq5/2$
\cite{Fielding2020}.  Their model connects the increase of cooling area to a
hot-gas inflow scaling close to $\tcool^{-1/4}$ in the strong-cooling regime.
The result explains why a planar interface area can severely underestimate the
total emission, but the fitted fractal dimension and cutoff scale should not be
treated as universal constants.

Tan, Oh, and Gronke organized the problem using turbulent-combustion theory
\cite{Tan2021}.  In terms of the outer-scale turbulent velocity $u'$ and size
$L$, their simulations distinguish
\begin{align}
 Q,\;v_{\rm in} &\propto {u'}^{1/2}(L/\tcool)^{1/2},
 &&\Da<1 \quad\text{(weak cooling)}, \label{eq:weakscaling}\\
 Q,\;v_{\rm in} &\propto {u'}^{3/4}(L/\tcool)^{1/4},
 &&\Da>1 \quad\text{(strong cooling)}, \label{eq:strongscaling}
\end{align}
where $Q=L_{\rm rad}/A_{\rm projected}$ is the bolometric surface brightness.
The apparently conflicting $\tcool^{-1/2}$ and $\tcool^{-1/4}$ results in the
literature can therefore describe different $\Da$ regimes.  These relations
predict a bolometric energy flux, not a wavelength-dependent spectrum.

\subsection{Local front structure and observable line predictions}

Tan and Oh made a further distinction that is especially important for the
present purpose \cite{TanOh2021}.  Global mass entrainment and bolometric
cooling can converge when the outer turbulent scale is resolved, even if the
microscopic interface is not.  In contrast, the temperature PDF, ion columns,
and line ratios depend on the local conductive--cooling front.  A steady
one-dimensional front obeys
\begin{equation}
 \frac{d}{dx}\left(\kappa\frac{dT}{dx}\right)
 =j_xc_p\frac{dT}{dx}+\rho\mathcal{L}(T),
 \label{eq:front}
\end{equation}
where $j_x=\rho v_x$ is the mass flux and
$\rho\mathcal{L}=n^2\Lambda-n\Gamma$.  The local temperature scale includes
the Field length
\begin{equation}
 \lambda_{\rm F}=\left(\frac{\kappa T}{n^2\Lambda}\right)^{1/2}.
\end{equation}
Their 1-D front model reproduced temperature distributions, column densities,
and line ratios from 3-D hydrodynamic calculations.  It also showed why the
temperature dependence of $\kappa$ matters: Spitzer
$\kappa\propto T^{5/2}$, constant conductivity, turbulent diffusion, and
grid-scale numerical diffusion weight the emitting temperatures differently.
This work supplies a practical subgrid model, but its published spectral tests
mostly assumed CIE; an RTML spectrum should combine this local-front treatment
with NEI whenever the time-scale test in Section~\ref{sec:nei} fails.

\subsection{High-Mach-number and magnetized extensions}

Yang and Ji extended plane-parallel RTML calculations to supersonic shear
\cite{YangJi2023}.  Below $\mathcal{M}\sim1$, their bolometric surface
brightness and intermediate-ion columns rose roughly as $\mathcal{M}^{1/2}$.
Above unity they saturated, because the velocity within the cooling/mixing
zone was limited by the local sound speed.  Supersonic layers developed a
two-zone structure: a nearly Mach-independent cooling/mixing zone and a
broader turbulent zone.  Turbulent dissipation, rather than hot-phase enthalpy
alone, became the dominant power source and cold gas tended toward evaporation
rather than condensation.

Recent MHD studies sharpen the magnetic-field caveat.  Zhao and Bai found that
even initially weak fields can be amplified enough to lower the surface
brightness and hot-gas inflow, through both magnetic pressure and suppression
of turbulent mixing \cite{ZhaoBai2023}.  Das and Gronke likewise found strong
suppression in idealized plane-parallel RTMLs, while driven turbulent boxes
showed smaller changes in net cold-gas growth because magnetic fields also
altered global morphology and contact area \cite{DasGronke2024}.  Therefore an
HD scaling cannot simply be multiplied by a universal magnetic correction:
the Alfv\'en Mach number, initial topology, anisotropic conduction, and global
geometry must be specified.

\subsection{Applications and present observational status}

In cloud--wind simulations, Gronke and Oh identified the useful survival and
growth condition
\begin{equation}
 t_{\rm cool,mix}\lesssim t_{\rm cc},
\end{equation}
where $t_{\rm cc}$ is the cloud-crushing time \cite{GronkeOh2018}.  Cooling of
mixed gas then creates new comoving cold material and can allow the cold mass
to exceed its initial value.  This is an application of RTML physics, although
the curved, evolving cloud surface is not identical to a periodic
plane-parallel layer.

Fielding and Bryan promoted this local cloud--wind exchange into a
steady-state model for an entire multiphase galactic wind
\cite{FieldingBryan2022}.  Their framework evolves the volume-filling hot
phase and an embedded cold-cloud population simultaneously, treating
cloud growth or destruction, drag, kinetic-energy thermalization, and metal
exchange as mass, momentum, and energy source terms for the hot flow.  It
does not resolve an RTML, but it supplies the missing population and geometry
closure needed to translate a single-interface calculation into a radial
wind luminosity.

For cosmological cold streams, Mandelker et al. found that the analogous ratio
of the mixed-gas cooling time to the shear-layer growth time controls whether a
stream is disrupted or gains mass \cite{Mandelker2020}.  Their follow-up
radiation model showed that KHI-induced dissipation could reach the luminosity
scale of observed Ly$\alpha$ blobs in massive high-redshift haloes
\cite{MandelkerLya2020}.  That comparison remains model-dependent because
Ly$\alpha$ requires self-shielding, resonant transfer, dust, and the external
UV field; bolometric RTML power cannot be equated directly to escaping
Ly$\alpha$ luminosity.

Taken together, existing studies support RTMLs as a physically plausible
source of intermediate temperature ions and interface radiation.  They do not
yet establish a unique spectral template.  Per-interface columns are often
below observed CGM columns, with some models requiring roughly
$10^2$--$10^3$ interfaces along a sightline \cite{TanOh2021,Ji2019}.  The
number and orientation of interfaces, NEI history, magnetic topology, and
radiative transfer are therefore part of the observable model rather than
nuisance details.

\section{What must be calculated for a spectrum}

\subsection{From hydrodynamics to an emissivity distribution}

For an optically thin cell $k$, the spectral luminosity is
\begin{equation}
 L_\lambda=\sum_k \epsilon_\lambda(T_k,\nelec{}_{k},\boldsymbol{f}_k,
 Z, J_\nu)\,V_k,
 \label{eq:sum}
\end{equation}
where $\boldsymbol{f}$ is the vector of ion fractions and $J_\nu$ is any
external or internally generated radiation field.  It is often useful to
store a differential emission measure (DEM),
\begin{equation}
 \frac{d\mathrm{EM}}{d\log T}=\sum_{k\in d\log T}\nelec{}_{k}\nH{}_{k}V_k,
 \qquad \mathrm{EM}=\int\nelec\nH\,dV,
 \label{eq:dem}
\end{equation}
and convolve it with a CIE contribution function $G_\lambda(T)$:
\begin{equation}
 L_\lambda^{\rm CIE}=\int G_\lambda(T;\nelec,Z)\,
 \frac{d\mathrm{EM}}{d\log T}\,d\log T.
 \label{eq:demcie}
\end{equation}
Equation~(\ref{eq:demcie}) is fast and excellent for a CIE reference model,
but it is invalid if $\boldsymbol{f}$ is not the CIE fraction at the local
temperature.  Do not use a bolometric $\Lambda(T)$ as $G_\lambda$; the former
has already summed over all radiative channels.

For a resolved transition $u\rightarrow l$ the volume emissivity is
\begin{equation}
 \epsilon_{ul}=n_u A_{ul}h\nu_{ul}.
 \label{eq:lineem}
\end{equation}
At low density, statistical equilibrium for the excited levels can usually be
solved at fixed ion fraction; collisional excitation, radiative decay,
recombination cascades, and (when relevant) photoexcitation enter the rate
matrix.  The continuum must be added on the same wavelength grid: free--free,
free--bound, two-photon, and, where appropriate, dust emission.  CHIANTI is
well suited to optically thin UV/EUV/optical line and continuum synthesis;
AtomDB/APEC is a standard choice for CIE X-ray spectra \cite{CHIANTI,AtomDB}.

\subsection{Why NEI is normally necessary in RTMLs}
\label{sec:nei}

For each element $X$ and charge state $q$, evolve
\begin{align}
 \frac{df_{X,q}}{dt}={}&\nelec\big[f_{X,q-1}C_{q-1}(T)+f_{X,q+1}\alpha_{q+1}(T)\\
 &-f_{X,q}\{C_q(T)+\alpha_q(T)\}\big]+\mathcal{P}_{X,q}(J_\nu)
 +\mathcal{X}_{X,q},
 \label{eq:nei}
\end{align}
with $\sum_qf_{X,q}=1$.  $C$, $\alpha$, $\mathcal{P}$, and $\mathcal{X}$
denote collisional ionization, total recombination, photoionization, and other
processes (e.g. charge exchange).  Couple this stiff system to the energy
equation because the emissivity and cooling rate depend on the evolving
fractions.  A simple decision rule is to calculate CIE only if both
$t_{\rm ion}$ and $t_{\rm rec}$ are short compared with \emph{both}
$\tcool$ and the local mixing/advection time.  Otherwise use NEI.

This is not a small correction in a mixing layer: mixing can be faster than
ionization, and cooling can be faster than recombination.  Kwak and Shelton
and Ji, Oh, and Masterson find enhanced high-ion columns in NEI relative to
CIE in relevant RTML regimes \cite{Kwak2010,Ji2019}.  A spectrum intended for
O\,VI, N\,V, C\,IV, Si\,IV, or soft-X-ray diagnostics should therefore report
whether it is CIE, NEI, or photoionization+NEI.

\section{Recommended calculation workflow}

\begin{enumerate}[leftmargin=*,label=\textbf{Step \arabic*.}]
\item \textbf{Specify the physical experiment.}  Record $T_c,T_h$, pressure
or densities, relative velocity, metallicity and abundance pattern, magnetic
field/conduction assumptions, external radiation field, and the target
viewing geometry.  State whether the calculation is isochoric, isobaric, or
fully hydrodynamic.

\item \textbf{Obtain a converged flow history.}  Run a
hydrodynamic/MHD calculation with a cooling treatment consistent with the
chosen abundances and radiation field, or make a one-dimensional cooling-flow
surrogate.  Save $(\rho,T,\nelec,\boldsymbol{v})$ and either Lagrangian tracer
histories or sufficiently frequent snapshots.  Check that total radiative
losses converge with resolution and close the energy budget in
Equation~(\ref{eq:energy}).

\item \textbf{Solve the ionization history.}  Prefer NEI integrated along
tracers (or together with the flow).  Include H, He, C, N, O, Ne, Mg, Si, S,
and Fe at minimum for broad UV--X-ray coverage, plus charge exchange near the
cool interface.  Use an implicit/BDF solver or another positivity-preserving
stiff integrator.  As a controlled diagnostic, generate an otherwise identical
CIE spectrum and quote the ratio line by line.

\item \textbf{Evaluate atomic emissivities.}  For each saved cell/tracer,
compute level populations and $\epsilon_\lambda$ using its $T,\nelec$ and NEI
fractions.  CHIANTI/ChiantiPy is convenient for optically thin line work;
AtomDB/PyAtomDB is convenient for CIE X-rays.  For photoionized, optically
thick, or radiation-coupled zones, use a spectral-synthesis/transfer code such
as \textsc{Cloudy}, but supply the non-equilibrium state or a sequence of
short history steps rather than silently imposing equilibrium.  The local
\textsc{Mappings~V} installation is an additional, particularly useful route:
precompute cooling, ionization, and emissivity tables with the same abundance
and radiation assumptions used by the dynamics, or evolve representative
tracer histories through its time-dependent single-zone solver.

\item \textbf{Project and transfer the radiation.}  The intrinsic surface
brightness along a ray is
\begin{equation}
 I_\lambda=\frac{1}{4\pi}\int\epsilon_\lambda(s)
 \exp[-\tau_\lambda(s)]\,ds.
 \label{eq:transfer}
\end{equation}
For an optically thin forbidden line set $\tau_\lambda=0$.  Treat resonance
lines (e.g. Ly$\alpha$, C\,IV, O\,VI), dust UV extinction, and X-ray
photoelectric absorption explicitly whenever their columns make them
non-negligible.  Add thermal and turbulent Doppler shifts before binning;
then convolve with the telescope line-spread function and effective area.

\item \textbf{Validate and publish reproducibly.}  Compare the integrated
model spectrum against the cooling loss, $\int L_\lambda d\lambda\simeq
L_{\rm rad}$ after accounting for untracked bands.  Test atomic-data version,
time cadence, resolution, abundances, CIE versus NEI, and viewing angle.
Publish line luminosities, intrinsic and attenuated spectra, and the DEM,
rather than only selected line ratios.
\end{enumerate}

\section{Code assessment and practical implementation}
\label{sec:codes}

\subsection{Selection criteria}

The relevant question is not which code has the largest collection of
features, but which one minimizes the amount of new microphysics required for
the intended observable.  A practical RTML code should ideally provide:
\begin{enumerate}[leftmargin=*]
\item multidimensional HD/MHD with a controllable plane-parallel shear layer;
\item passive-scalar or species advection and a stiff NEI integrator;
\item all charge states needed for C\,IV, N\,V, and O\,VI;
\item radiative losses evaluated from the same instantaneous ion fractions
that are advected by the flow;
\item optional thermal conduction and a resolvable Field length; and
\item an export or post-processing route from cell states to line and
continuum emissivities.
\end{enumerate}
No candidate considered here supplies unrestricted three-dimensional dynamics,
full NEI, arbitrary-geometry frequency-dependent transfer, and an
observation-ready spectrum without some additional work.

\subsection{Comparison of the main candidates}

\begin{longtable}{
 >{\raggedright\arraybackslash}p{2.4cm}
 >{\raggedright\arraybackslash}p{3.4cm}
 >{\raggedright\arraybackslash}p{4.5cm}
 >{\raggedright\arraybackslash}p{3.4cm}}
\toprule
\textbf{Code} & \textbf{Existing strengths} &
\textbf{Missing or limiting feature} & \textbf{Recommended role}\\
\midrule
\endfirsthead
\toprule
\textbf{Code} & \textbf{Existing strengths} &
\textbf{Missing or limiting feature} & \textbf{Recommended role}\\
\midrule
\endhead
\textsc{Pluto} with the Sarkar--Sternberg--Gnat extension &
HD/MHD; full ion stages of H, He, C, N, O, Ne, Mg, Si, S, and Fe;
ion-by-ion cooling based on \textsc{Cloudy}-17; charge exchange and
photoionization; public source and an NEI spectral post-processor &
The published frequency-dependent transfer solver assumes spherical
symmetry; the extension is based on \textsc{Pluto} 4.0 and therefore needs
build and regression testing before a production 3-D run &
\textbf{Most direct public starting point} for a plane-parallel NEI RTML;
use its 3-D ion network and optically thin spectral post-processing, not its
spherical transfer solver \cite{Sarkar2021}\\
\textsc{Maihem} &
Modified \textsc{Flash} with NEI, ion-by-ion cooling, photoionization and
photoheating; 84 species across 13 elements; demonstrated non-CIE line
synthesis with frozen ion fractions passed to \textsc{Cloudy} &
No clearly maintained general public distribution was identified; obtaining
and validating the exact source version may require collaboration with the
authors &
\textbf{Best integrated physical match if the source is available}
\cite{Gray2019}\\
\textsc{Flash} 4.8 &
Mature AMR HD/MHD; an \texttt{Ionize/Nei} unit that advects every charge
state for a selectable set including C, N, and O; explicit/implicit isotropic
and anisotropic conduction &
The standard NEI unit does not feed ionization/recombination losses back into
the energy equation and uses the Summers (1974) rate file; it does not
produce a spectrum automatically &
\textbf{Most robust fallback}: add modern ion-dependent cooling and a
\textsc{Chianti}/\textsc{Cloudy} emissivity layer
\cite{Fryxell2000,FLASHGuide}\\
\textsc{Athena++} &
Modern, efficient AMR MHD framework; already used for plane-parallel RTML
calculations with cooling and conduction &
Published RTML setups usually use CIE cooling and synthetic spectra rather
than an in-flow full-metal NEI network &
Excellent CIE dynamics baseline or a platform for a new custom NEI module,
but not the shortest route to NEI emission \cite{Stone2020,TanOh2021}\\
Standard \textsc{Pluto} MINEq &
NEI source integration and self-consistent optically thin cooling from 29
species; designed for emission-line diagnostics &
The older network includes only the first five stages of C, N, O, Ne, and S,
so it does not by itself contain O\,VI; superseded for this project by the
full-ion extension &
Useful validation case, not the preferred production network
\cite{Tesileanu2008}\\
\textsc{Mappings~V} 5.2.1 (local checkout) &
Thirty-element atomic setup; fixed-temperature CIE, time-dependent
ionization, coupled time-dependent cooling and recombination, isochoric or
isobaric evolution, detailed line/continuum emissivities, and radiative-shock
models with self-consistent precursor calculations &
It is a standalone, mostly interactive Fortran microphysics/spectral code,
not a multidimensional turbulence or MHD solver; arbitrary 3-D RTML histories
require a wrapper or precomputed tables, and one invocation per simulation
cell would be impractical &
\textbf{Preferred local microphysics and validation engine}: generate
CIE/NEI cooling and emissivity tables, post-process representative tracer
histories, and benchmark the atomic physics independently of the hydro code
\cite{Sutherland2017}\\
{\footnotesize\shortstack[l]{\textsc{Multiphase}\\
\textsc{GalacticWind}}} &
Compact Python implementation of the Fielding--Bryan steady multiphase-wind
model; bidirectional hot/cloud mass, momentum, energy, and metal exchange;
RTML-informed cloud growth and ablation; radial cloud area and number-density
closure &
It neither resolves KH turbulence nor evolves NEI ions, line emissivities,
MHD, conduction, or transfer; the example evolves one characteristic cloud
population and uses equilibrium Wiersma et al.\ cooling tables &
\textbf{Preferred galaxy-scale population wrapper}: convert a calibrated
single-interface RTML flux into radial luminosity, surface-brightness, and
velocity-profile predictions for a multiphase wind
\cite{FieldingBryan2022,Wiersma2009}\\
\textsc{Torus} (local checkout) &
1-D/2-D/3-D AMR, Eulerian hydrodynamics, Monte Carlo radiation transport,
photoionization equilibrium, dust, and line-transfer machinery; a
\texttt{geometry="kelvin"} hydrodynamic initial condition is present &
No ready-made collisionally ionized full C/N/O NEI network, self-consistent
RTML metal cooling, MHD, or conductive front model was found; C\,IV is
present in the local ion list, whereas N\,V and the higher oxygen stages
needed for O\,VI are disabled or absent &
Use as a \textbf{transfer/synthetic-observation post-processor} after importing
the flow and NEI ion fractions, especially when dust, absorption, scattering,
or complex viewing geometry matters \cite{Harries2019}\\
\bottomrule
\end{longtable}

\subsection{Role of the local \textsc{Mappings~V} checkout}
\label{sec:mappings}

The inspected source tree is the public \textsc{Mappings~V} version 5.2.1.
Its full preference set loads 30 elements, including all carbon, nitrogen,
and oxygen ion stages needed for C\,IV, N\,V, and O\,VI.  The single-zone
driver provides equilibrium ionization at fixed temperature, equilibrium
ionization and temperature, time-dependent ionization at either fixed or
equilibrium temperature, and coupled time-dependent ionization and
temperature.  For the last case the density evolution may be chosen as
isochoric or isobaric.  The supplied cooling examples include both CIE and
NEI cooling trajectories, while the cooling and spectrum routines can write
element-resolved cooling, ionization fractions, line emissivities, and local
or integrated spectra.  These are exactly the ingredients required to
construct controlled RTML spectral benchmarks.

The radiative-shock module adds a second useful test.  It advances the
post-shock ionization and cooling zones using atomic, cooling, and flow time
scales and iterates the upstream photoionized precursor
\cite{Sutherland2017}.  Such a shock spectrum must \emph{not} be identified
with a shear-driven RTML spectrum: compression, velocity history, and
self-irradiation differ.  It is nevertheless a stringent test of the same
ionization, cooling, and emissivity machinery over the temperature range
occupied by a radiative mixing layer.

There are four recommended uses:
\begin{enumerate}[leftmargin=*]
\item generate a version-pinned CIE cooling curve and CIE emissivity grid for
the initial hydro and spectrum baseline;
\item generate isochoric and isobaric NEI cooling sequences and the associated
C\,IV/N\,V/O\,VI contribution functions;
\item post-process a representative set of Lagrangian thermal histories from
the RTML simulation; and
\item compare integrated cooling and selected line luminosities against
\textsc{Chianti}, \textsc{Cloudy}, and the \textsc{Pluto} full-ion
post-processor to expose atomic-data or normalization differences.
\end{enumerate}

For CIE, a table indexed by $(T,\nelec,Z,J_\nu)$ can be interpolated cell by
cell.  The analogous NEI table cannot generally be indexed by only the
\emph{current} $(T,\nelec)$, because two parcels at the same temperature can
retain different ion fractions after different mixing and cooling histories.
The production calculation must therefore do one of the following: evolve
the ions inside the hydro code; save tracer histories and pass them through
\textsc{Mappings}; or construct a deliberately restricted library of
isochoric/isobaric histories and quantify the approximation.  A final
snapshot containing only density and temperature is insufficient for a
general NEI spectrum.

Directly launching the interactive \textsc{Mappings} executable once for
every cell and output time would also be prohibitively expensive.  The
practical interface is a scripted batch wrapper that first computes tables
on a grid of physical conditions and then uses interpolation during
post-processing.  If the hydro cooling is taken from those tables, the
spectral calculation must retain the same abundances, dust assumptions,
radiation field, and atomic-data version.  A separate background metal
cooling term must not be added on top of the \textsc{Mappings} loss, because
that would double count radiative channels.

At the time of inspection the local source tree contained no built
\texttt{map52} executable.  Its documented build target is
\texttt{make build}, with a default installation under
\texttt{\string~/mappings520}.  Before scientific use, build it in an
isolated test location, run the supplied CIE and NEI cooling examples, and
archive the exact preference, abundance, dust, and atomic-data files with
the generated tables.

\subsection{From a local RTML to a multiphase galactic wind}
\label{sec:fieldingbryan}

The public \textsc{MultiphaseGalacticWind} repository is a compact
implementation of the Fielding--Bryan analytic framework
\cite{FieldingBryan2022}.  It first constructs a Chevalier--Clegg-like hot
wind and then solves radial ordinary differential equations for
\[
 \{v_{\rm w},\rho_{\rm w},P_{\rm w},\rho_{\rm w}Z_{\rm w},
   M_{\rm c},v_{\rm c},Z_{\rm c}\}.
\]
The hot and cold phases exchange mass, momentum, energy, and metals.  The
cloud is pressure confined; its acceleration includes ram-pressure drag,
momentum gained with condensing material, and gravity, while the hot wind
responds to all of these source terms.  This bidirectional coupling is more
useful for a global wind than treating clouds as passive test particles.

The RTML closure enters through
\begin{equation}
 T_{\rm mix}=(T_{\rm w}T_{\rm c})^{1/2},\qquad
 Z_{\rm mix}=(Z_{\rm w}Z_{\rm c})^{1/2},\qquad
 \xi=\frac{r_{\rm c}}{v_{\rm turb}t_{\rm cool,mix}}.
 \label{eq:windxi}
\end{equation}
The example implementation uses a cloud-growth factor proportional to
$\xi^{1/2}$ for $\xi<1$ and $\xi^{1/4}$ for $\xi>1$, corresponding to the
weak- and strong-cooling RTML scalings, together with a separately
parameterized turbulent ablation rate and an effective cooling-area boost.
It therefore provides a concrete route from the local scalings in
Equations~(\ref{eq:weakscaling}) and (\ref{eq:strongscaling}) to a radial
mass-loading model.  The coefficients and density-contrast powers remain
closures calibrated by idealized simulations and must be varied rather than
treated as universal constants.

Its most important contribution to spectral modeling is geometrical.  For a
steady cloud number flux $\dot N_{\rm c}$ in solid angle $\Omega$, the cloud
number density is
\begin{equation}
 n_{\rm c}(r)=\frac{\dot N_{\rm c}(r)}
 {\Omega r^2v_{\rm c}(r)}.
 \label{eq:cloudnumber}
\end{equation}
If a local RTML calculation supplies a spectral surface flux
$F_\lambda^{\rm RTML}$ and the effective emitting area of one cloud is
$A_{\rm eff}$, then
\begin{align}
 L_\lambda^{\rm c}(r)&=A_{\rm eff}(r)
 F_\lambda^{\rm RTML}
 [P,T_{\rm w},T_{\rm c},v_{\rm rel},Z,\mathrm{history}],\\
 \frac{dL_\lambda}{dr}
 &=\Omega r^2 n_{\rm c}L_\lambda^{\rm c}
 =\frac{\dot N_{\rm c}}{v_{\rm c}}L_\lambda^{\rm c}.
 \label{eq:windlum}
\end{align}
This converts a planar or single-cloud emissivity calculation into a
galaxy-scale radial luminosity.  Doppler shifting each contribution by the
projected $v_{\rm c}(r)$ and $v_{\rm w}(r)$ then gives an intrinsic line
profile before transfer.  The needed
$F_\lambda^{\rm RTML}$ should come from \textsc{Mappings} tables or resolved
\textsc{Pluto}/\textsc{Flash} interface calculations, not from the
bolometric cooling curve alone.

The repository's bundled cooling table is based on the redshift-zero
Wiersma et al.\ UV-background calculation \cite{Wiersma2009}.  It is useful
for the steady wind energy equation but supplies neither the NEI state nor
C\,IV/N\,V/O\,VI line emissivities.  The code also assumes steady radial
flow and a representative cloud rather than a time-dependent distribution of
cloud sizes and histories.  Accordingly, it should be used as a
\emph{population wrapper} around RTML microphysics, not as a replacement for
the hydro, NEI, spectral, or transfer stages.

\subsection{Revised practical recommendation}

For work that can start from publicly identified source, the most economical
choice is the Sarkar--Sternberg--Gnat \textsc{Pluto} extension
\cite{Sarkar2021}.  Its ionization network is valid over
$5\times10^3$--$3\times10^8\,\mathrm{K}$, includes all stages needed for
C\,IV/N\,V/O\,VI, calculates cooling for each ion--electron pair on the fly,
and supplies a post-processing tool for projected columns and spectra.  The
ion network is compatible with three-dimensional MHD even though its
frequency-dependent radiation-transfer module is restricted to spherical
symmetry.  A plane-parallel RTML should therefore use the 3-D MHD+NEI part and
perform optically thin spectral integration separately.  The local
\textsc{Mappings} code should then be used as an independent one-zone
microphysics benchmark and, where its atomic data are preferred, as a
secondary emissivity-table generator.

If the \textsc{Maihem} source can be obtained and built reproducibly, it is an
even closer physical match.  It evolves the ionization state and computes
cooling ion by ion, and published applications pass the non-equilibrium
fractions to \textsc{Cloudy} to calculate line emissivities.  This procedure
has produced strong C\,IV $\lambda1549$ and O\,VI $\lambda1037$ emission in
rapidly cooling flows \cite{Gray2019}.  Its practical ranking is limited by
source availability rather than physics.

\textsc{Flash} 4.8 remains the safest independently maintainable alternative.
Its built-in NEI unit already removes the largest bookkeeping burden---the
advection and reaction solution for all C/N/O stages---but a production RTML
calculation must replace or augment the old rate data and must couple the
instantaneous ion state to the energy source.  The documented defaults
$10^4<T<10^7\,\mathrm{K}$ and
$1<\nelec<10^{12}\,\mathrm{cm}^{-3}$ also require special attention:
the default lower electron-density bound is unsuitable for much of the CGM
unless it is changed and validated \cite{FLASHGuide}.

Consequently the recommended order is:
\begin{enumerate}[leftmargin=*]
\item \textbf{Shortest self-consistent NEI route:} the full-ion
\textsc{Pluto} extension for dynamics and ion evolution, with
\textsc{Mappings} and \textsc{Chianti} as independent spectral checks;
\item \textbf{most maintainable modular route:} \textsc{Flash} 4.8 for AMR
HD/MHD plus version-pinned \textsc{Mappings} cooling/emissivity tables and a
new coupling layer;
\item \textbf{best integrated physics if obtainable:} \textsc{Maihem}, again
benchmarked against \textsc{Mappings};
\item \textbf{fast CIE control calculation:} \textsc{Athena++} or standard
\textsc{Pluto} with a \textsc{Mappings}/\textsc{Chianti} tabulated CIE cooling
curve; and
\item \textbf{galaxy-scale population wrapper:}
\textsc{MultiphaseGalacticWind}, after calibrating its cloud-growth,
ablation, and emitting-area closures with the local RTML calculation; and
\item \textbf{radiation-transfer finishing stage:} \textsc{Torus}, only after
the required ion fractions and intrinsic emissivities have been supplied.
\end{enumerate}

Thus the most realistic near-term stack for the available local codes is
\[
\begin{aligned}
 \textsc{Pluto/Flash}\;(\text{RTML dynamics})
 &\longrightarrow \textsc{Mappings~V}\;(\text{NEI emission})\\
 &\longrightarrow \textsc{MultiphaseGalacticWind}\;
   (\text{optional population model})\\
 &\longrightarrow \textsc{Torus}\;(\text{optional transfer}).
\end{aligned}
\]
\textsc{MultiphaseGalacticWind} is omitted for a single-interface study.
\textsc{Torus} is omitted from the first production pass if all diagnostic
lines are demonstrably optically thin and foreground attenuation can be
applied analytically.

\subsection{A concrete production architecture}

The proposed calculation should be developed in five increasingly expensive
levels.

\paragraph{Level 1: hydrodynamic CIE benchmark.}
Initialize a pressure-balanced hot/cold interface in a Cartesian box with the
shear direction periodic, a transverse domain large enough to contain the
growing layer, and seeded perturbations whose spectrum is recorded.  Run 2-D
resolution and cooling-strength scans before the 3-D production model.
Compute the bolometric loss, DEM, mixing-layer width, entrainment velocity,
and CIE C\,IV/N\,V/O\,VI columns.  This stage tests Equations
(\ref{eq:energy})--(\ref{eq:strongscaling}) without conflating a hydro error
with an NEI error.  Generate the baseline cooling and emissivity tables with
\textsc{Mappings}, and cross-check a sparse set of temperatures with
\textsc{Chianti}.

\paragraph{Level 2: dynamically coupled NEI.}
Evolve at least H, He, C, N, and O as conserved species with reaction source
terms.  Other important coolants may either be evolved or included through a
carefully separated background cooling term.  The two contributions must not
double count the same metal lines.  Advance reactions and cooling with
subcycling or an implicit stiff solver, enforce positivity and elemental
normalization, and recompute $\nelec$ and the mean molecular weight from the
instantaneous ion state.  Compare the NEI and CIE cooling histories for
identical fluid initial conditions.  Validate the coupled solver against the
\textsc{Mappings} isochoric and isobaric single-zone cooling sequences before
running a turbulent box.

\paragraph{Level 3: intrinsic emission cube and spectrum.}
For each transition $u\rightarrow l$, calculate
\begin{equation}
 \epsilon_\nu(\boldsymbol{x})=
 \nelec n_{X,q}\,q_{ul}(T,\nelec)\,h\nu_{ul}\,
 \phi_\nu[T,\boldsymbol{v},v_{\rm turb}],
 \label{eq:profileem}
\end{equation}
where $\phi_\nu$ contains thermal broadening, the cell line-of-sight velocity,
and any explicitly defined unresolved velocity dispersion.  Obtain
$q_{ul}$ and continuum coefficients from a version-pinned atomic data set,
using either \textsc{Mappings} or \textsc{Chianti}; sparse one-zone
\textsc{Cloudy} calculations with frozen NEI fractions provide a useful
cross-check.  If \textsc{Mappings} is used, interpolate precomputed
emissivities rather than starting one interactive model per cell.  Produce
both a cell emissivity cube and the volume-integrated spectrum so that
inclination and aperture can be varied without rerunning the dynamics.

\paragraph{Level 4: optional wind-scale population synthesis.}
If the target is an integrated galactic wind rather than one interface,
interpolate the RTML surface-flux library along the radial solution of
Section~\ref{sec:fieldingbryan}.  Calculate $A_{\rm eff}$, cloud number
density, $dL_\lambda/dr$, projected velocity, and aperture dependence.
Vary the cloud injection history or mass spectrum and the uncertain
growth/ablation coefficients.  Energy radiated by the interfaces must remain
consistent with the mass and enthalpy exchange applied to the hot wind.

\paragraph{Level 5: transfer and instrument model.}
First calculate the optically thin result using
Equation~(\ref{eq:transfer}).  Only if the inferred line-center optical depth,
dust opacity, or photoelectric absorption is significant should the emission
cube be passed to \textsc{Torus} or another transfer solver.  Importing NEI
fractions into \textsc{Torus} will require a defined data interface and an
extension of its atomic model for N\,V and O\,VI; its current photoionization
equilibrium solution should not overwrite the imported NEI state.  Resonant
scattering must preserve photon frequency and direction information rather
than being represented by a scalar escape fraction when line profiles are an
observable.

\subsection{Minimum verification matrix}

Before interpreting a spectrum, vary at least the grid scale, cooling
normalization, conduction coefficient, magnetic-field strength/orientation,
and ionization model.  The minimum reported comparison is
\[
 \{\mathrm{CIE},\mathrm{NEI}\}\times
 \{\mathrm{HD},\mathrm{MHD}\}\times
 \{\kappa=0,\kappa=\kappa_{\rm adopted}\},
\]
with matched initial conditions.  Verify conservation of every elemental
abundance, positivity of all fractions, cooling-energy closure, convergence
of the DEM and line luminosities, and the reproduction of one-zone CIE and
isochoric/isobaric NEI benchmarks.  Because emissivity is approximately
quadratic in density, agreement in total mass or column density alone is not
a sufficient emission test.

\section{A minimal, useful surrogate: a steady isobaric cooling flow}

When a full 3-D RTML calculation is unavailable, a steady cooling-flow model
gives a transparent first estimate.  Assume a mass flux per interface area
$\dot m/A$ and evolve a parcel from its initial mixed temperature under
constant pressure.  With enthalpy per mass $h$, solve
\begin{equation}
 \frac{dh}{dt}=-\frac{\nelec\nH\Lambda_{\rm net}(T,\boldsymbol f)}{\rho},
 \quad \frac{d\boldsymbol f}{dt}=\boldsymbol R(T,\nelec,J_\nu)\boldsymbol f,
 \label{eq:parcel}
\end{equation}
then compute a line luminosity per interface area,
\begin{equation}
 \frac{L_{ul}}{A}=\frac{\dot m}{A}\int \frac{\epsilon_{ul}}{\rho}\,dt.
 \label{eq:coolflowline}
\end{equation}
This has two benefits: it makes the assumed energy/mass flux explicit, and it
handles ionization lag.  Its limitation is equally explicit: it cannot predict
the entrainment rate, corrugated area, or velocity-dependent line profile.
Those require a turbulence model or simulation.  Use it as a benchmark and
parameter study, not as a substitute for resolved RTML dynamics.  The
time-dependent, isobaric option in the local \textsc{Mappings} single-zone
solver is a direct implementation path for this surrogate, while the
isochoric option brackets sensitivity to the thermodynamic closure.

\section{Diagnostics and observables}

The expected output is multi-band.  The $\sim10^{4.5}$--$10^{5.5}$ K material
can produce UV high-ion diagnostics (C\,IV, Si\,IV, N\,V, O\,VI); hotter gas
contributes EUV/soft-X-ray lines and continuum; recombined material can emit
optical/IR lines if it remains ionized or is externally illuminated.  The
classical mixing-layer models of Slavin, Shull, and Begelman explicitly
connected high ions and UV/optical emission to a turbulent interface
\cite{Slavin1993}.

Interpret line ratios jointly with the DEM and absorption columns:
\begin{equation}
 N_{X,q}=\int n_X f_{X,q}\,ds.
\end{equation}
Emission scales approximately as density squared in the optically thin,
low-density limit, while absorption scales linearly with density.  Combining
them constrains filling factor and path length far better than either alone.
NEI may make an ion appear at a temperature where its CIE contribution
function is small; a ``formation temperature'' is therefore a model-dependent
description, not a direct thermometer in an RTML.

\section{Numerical and physical failure modes}

\begin{itemize}[leftmargin=*]
\item \textbf{Cooling--spectrum mismatch:} a tabulated CIE cooling curve is
used in the dynamics while post-processing assumes NEI or a different
abundance set.  Recompute or quantify the inconsistency.
\item \textbf{Under-resolved numerical mixing:} numerical diffusion can set
the DEM and artificially boost cooling.  Demonstrate convergence of integrated
luminosity, DEM, and selected line luminosities, not merely morphology.
\item \textbf{Equilibrium by assumption:} CIE can miss the lagging high ions
in cooling/mixing gas.  Carry NEI fractions, at least in post-processing.
\item \textbf{Ignoring transfer:} optically thin synthesis is not adequate
for strong resonance lines or dusty/high-column sightlines.
\item \textbf{Comparing unlike quantities:} distinguish intrinsic luminosity,
surface brightness, attenuated flux, and detector counts; each has different
geometry and units.
\end{itemize}

\section{Literature context}

\begin{longtable}{p{3.1cm}p{3.2cm}p{7.0cm}}
\toprule
\textbf{Study} & \textbf{Main method} & \textbf{Reliable takeaway for spectral work}\\
\midrule
\endfirsthead
\toprule
\textbf{Study} & \textbf{Main method} & \textbf{Reliable takeaway for spectral work}\\
\midrule
\endhead
Begelman \& Fabian (1990) & Analytic energy balance & Establishes the enthalpy--cooling balance that remains a necessary global check.\\
Slavin, Shull \& Begelman (1993) & Turbulent-interface/NEI spectral models & Direct precedent for predicting high-ion emission and absorption from mixing layers.\\
Esquivel et al. (2006) & 3-D MHD + CIE cooling & Demonstrates time-dependent layer growth and magnetic sensitivity; CIE limits its spectral interpretation.\\
Kwak \& Shelton (2010) & Hydro + NEI & Mixing and cooling histories enhance high ions relative to CIE; time dependence matters.\\
Ji, Oh \& Masterson (2019) & 3-D HD/MHD + NEI + radiation & Simple $h\sim v_{\rm turb}\tcool$ scalings can fail; NEI changes high-ion columns by factors of a few.\\
Fielding et al. (2020) & 3-D hydro, fractal cooling front & The cooling surface is corrugated; use global energetics and resolution tests rather than planar geometry alone.\\
Tan, Oh \& Gronke (2021) & Turbulent-combustion analogy & $\Da$ organizes weak- and strong-cooling regimes and informs convergence tests.\\
Tan \& Oh (2021) & 1-D front + 3-D spectral tests & Global entrainment and line observables have different resolution requirements; local conduction controls the temperature PDF.\\
Yang \& Ji (2023) & High-Mach 3-D hydro & Surface brightness and ion columns saturate above $\mathcal{M}\sim1$; dissipation can drive evaporation.\\
Zhao \& Bai (2023); Das \& Gronke (2024) & 3-D MHD & Amplified weak fields suppress local mixing, while global morphology can partly compensate in turbulent systems.\\
Fielding \& Bryan (2022) & Steady two-phase wind ODEs with RTML-informed exchange & Converts local cloud growth/ablation into bidirectional mass, momentum, energy, and metal source terms; useful for interface population and radial luminosity modeling, not for NEI microphysics.\\
\bottomrule
\end{longtable}

\section{Conclusions}

An RTML radiation spectrum is a state-resolved emissivity calculation,
not a direct plot of $\Lambda(T)$.  The recommended deliverable is a
NEI-based intrinsic spectrum plus transferred spectra for specified sightlines,
accompanied by a DEM, ion columns, and an energy-closure test.  A CIE/DEM
spectrum remains valuable as a fast control model; differences from NEI then
become an interpretable prediction of the mixing and cooling history.

For implementation, the current first choice is the publicly available
full-ion NEI extension of \textsc{Pluto}, because it already connects
multidimensional MHD, the required C/N/O charge states, ion-dependent cooling,
and an emission-spectrum post-processor.  The local \textsc{Mappings~V} 5.2.1
checkout is the preferred complementary microphysics tool: it can supply CIE
and NEI cooling/emissivity tables, one-zone cooling-flow calculations, and
radiative-shock benchmarks, but it cannot form the multidimensional mixing
layer.  \textsc{Maihem} is the most integrated alternative if its source can
be obtained, and \textsc{Flash} 4.8 is the preferred robust platform when the
cooling/emissivity coupling will be developed locally.  The resulting
practical architecture is \textsc{Pluto}/\textsc{Flash} for dynamics,
\textsc{Mappings} for microphysics and intrinsic radiation,
\textsc{MultiphaseGalacticWind} when a local interface must be embedded in a
radial cloud population, and \textsc{Torus} only as the final transfer stage
for an externally computed NEI emission cube.

\begin{thebibliography}{99}
\footnotesize
\bibitem{Begelman1990} Begelman, M. C. \& Fabian, A. C. 1990, MNRAS, 244, 26P, doi:\href{https://doi.org/10.1093/mnras/244.1.26P}{10.1093/mnras/244.1.26P}.
\bibitem{Slavin1993} Slavin, J. D., Shull, J. M. \& Begelman, M. C. 1993, ApJ, 407, 83, doi:\href{https://doi.org/10.1086/172494}{10.1086/172494}.
\bibitem{Esquivel2006} Esquivel, A., Benjamin, R. A., Lazarian, A., Cho, J. \& Leitner, S. N. 2006, ApJ, 648, 1043, doi:\href{https://doi.org/10.1086/505338}{10.1086/505338}.
\bibitem{Kwak2010} Kwak, K. \& Shelton, R. L. 2010, ApJ, 719, 523, doi:\href{https://doi.org/10.1088/0004-637X/719/1/523}{10.1088/0004-637X/719/1/523}.
\bibitem{Gnat2007} Gnat, O. \& Sternberg, A. 2007, ApJS, 168, 213, doi:\href{https://doi.org/10.1086/509786}{10.1086/509786}.
\bibitem{Ji2019} Ji, S., Oh, S. P. \& Masterson, P. 2019, MNRAS, 487, 737, doi:\href{https://doi.org/10.1093/mnras/stz1248}{10.1093/mnras/stz1248}.
\bibitem{Fielding2020} Fielding, D. B., Ostriker, E. C., Bryan, G. L. \& Jermyn, A. S. 2020, ApJL, 894, L24, doi:\href{https://doi.org/10.3847/2041-8213/ab8d2c}{10.3847/2041-8213/ab8d2c}.
\bibitem{Tan2021} Tan, B., Oh, S. P. \& Gronke, M. 2021, MNRAS, 502, 3179, doi:\href{https://doi.org/10.1093/mnras/stab053}{10.1093/mnras/stab053}.
\bibitem{TanOh2021} Tan, B. \& Oh, S. P. 2021, MNRAS Letters, 508, L37, doi:\href{https://doi.org/10.1093/mnrasl/slab100}{10.1093/mnrasl/slab100}.
\bibitem{YangJi2023} Yang, Y. \& Ji, S. 2023, MNRAS, 520, 2148, doi:\href{https://doi.org/10.1093/mnras/stad264}{10.1093/mnras/stad264}.
\bibitem{ZhaoBai2023} Zhao, X. \& Bai, X.-N. 2023, MNRAS, 526, 4245, doi:\href{https://doi.org/10.1093/mnras/stad3011}{10.1093/mnras/stad3011}.
\bibitem{DasGronke2024} Das, H. K. \& Gronke, M. 2024, MNRAS, 527, 991, doi:\href{https://doi.org/10.1093/mnras/stad3125}{10.1093/mnras/stad3125}.
\bibitem{GronkeOh2018} Gronke, M. \& Oh, S. P. 2018, MNRAS Letters, 480, L111, doi:\href{https://doi.org/10.1093/mnrasl/sly131}{10.1093/mnrasl/sly131}.
\bibitem{Mandelker2020} Mandelker, N., Nagai, D., Aung, H., Dekel, A., Birnboim, Y. \& van den Bosch, F. C. 2020, MNRAS, 494, 2641, doi:\href{https://doi.org/10.1093/mnras/staa812}{10.1093/mnras/staa812}.
\bibitem{MandelkerLya2020} Mandelker, N., van den Bosch, F. C., Nagai, D., Dekel, A., Birnboim, Y. \& Aung, H. 2020, MNRAS, 498, 2415, doi:\href{https://doi.org/10.1093/mnras/staa2421}{10.1093/mnras/staa2421}.
\bibitem{CHIANTI} Del Zanna, G., et al. 2021, ApJ, 909, 38, doi:\href{https://doi.org/10.3847/1538-4357/abd8ce}{10.3847/1538-4357/abd8ce}; \href{https://www.chiantidatabase.org/}{CHIANTI database}.
\bibitem{AtomDB} Foster, A. R., et al. 2012, ApJ, 756, 128, doi:\href{https://doi.org/10.1088/0004-637X/756/2/128}{10.1088/0004-637X/756/2/128}; \href{https://www.atomdb.org/}{AtomDB}.
\bibitem{Fryxell2000} Fryxell, B., et al. 2000, ApJS, 131, 273, doi:\href{https://doi.org/10.1086/317361}{10.1086/317361}.
\bibitem{FLASHGuide} Flash Center for Computational Science 2024, \textit{FLASH User's Guide, Version 4.8}, \href{https://flash.rochester.edu/site/flashcode/user_support/flash_ug_devel/flash4_ug.html}{online manual}, especially the \href{https://flash.rochester.edu/site/flashcode/user_support/flash_ug_devel/node120.html}{Ionization Unit} and \href{https://flash.rochester.edu/site/flashcode/user_support/flash_ug_devel/node128.html}{Diffuse Unit}.
\bibitem{Tesileanu2008} Te\c{s}ileanu, O., Mignone, A. \& Massaglia, S. 2008, A\&A, 488, 429, doi:\href{https://doi.org/10.1051/0004-6361:200809461}{10.1051/0004-6361:200809461}.
\bibitem{Sarkar2021} Sarkar, K. C., Sternberg, A. \& Gnat, O. 2021, MNRAS, 503, 5807, doi:\href{https://doi.org/10.1093/mnras/stab578}{10.1093/mnras/stab578}; \href{https://gitlab.com/kartickchsarkar/pluto-neq-radiation}{public source repository}.
\bibitem{Sutherland2017} Sutherland, R. S. \& Dopita, M. A. 2017, ApJS, 229, 34, doi:\href{https://doi.org/10.3847/1538-4365/aa6541}{10.3847/1538-4365/aa6541}; \href{https://mappings.anu.edu.au/}{MAPPINGS public distribution}.
\bibitem{FieldingBryan2022} Fielding, D. B. \& Bryan, G. L. 2022, ApJ, 924, 82, doi:\href{https://doi.org/10.3847/1538-4357/ac2f41}{10.3847/1538-4357/ac2f41}; \href{https://github.com/dfielding14/MultiphaseGalacticWind/}{public example code}.
\bibitem{Wiersma2009} Wiersma, R. P. C., Schaye, J. \& Smith, B. D. 2009, MNRAS, 393, 99, doi:\href{https://doi.org/10.1111/j.1365-2966.2008.14191.x}{10.1111/j.1365-2966.2008.14191.x}.
\bibitem{Gray2019} Gray, W. J., Oey, M. S., Silich, S. \& Scannapieco, E. 2019, ApJ, 887, 161, doi:\href{https://doi.org/10.3847/1538-4357/ab510d}{10.3847/1538-4357/ab510d}.
\bibitem{Stone2020} Stone, J. M., Tomida, K., White, C. J. \& Felker, K. G. 2020, ApJS, 249, 4, doi:\href{https://doi.org/10.3847/1538-4365/ab929b}{10.3847/1538-4365/ab929b}.
\bibitem{Harries2019} Harries, T. J., Haworth, T. J., Acreman, D., Ali, A. \& Douglas, T. 2019, Astronomy and Computing, 27, 63, doi:\href{https://doi.org/10.1016/j.ascom.2019.03.002}{10.1016/j.ascom.2019.03.002}.
\end{thebibliography}

\end{document}
